\section{A uniform bound on natural-module composition length}
\label{pkt:section}

We continue in the setting of Proposition~\ref{grw:remaining-paths}.
The endpoint types have a common characteristic, field parameter
$q=p^f$, and classical family.  Write $n_i$ for their natural
dimensions, using the symplectic realization for a group of type
$B$ in characteristic two.  The natural scalar field is
$\mathbb F_q$, except in the unitary case, where it is
$\mathbb F_{q^2}$.  All ranks in this section may be assumed to be
at least nine.

For a comparison $i<j$, let $P=P_{ij}$ be the inner projection
in the upper simple group $T_j$, and let $M=M_{ij}$ be the
corresponding lower coordinate stabilizer in
$\operatorname{Aut}(T_j)$.  Proposition~\ref{lift:edge-lifts} supplies a
quasisimple component $E=E_{ij}$, normal in both $P$ and $M$, with
$P/E$ soluble.  In matrix representations we take the indicated
central covers.  The same stabilizer normalizes the intermediate
component on either edge of a shortcut.  These are structural
conclusions about the remaining edges, not consequences of the
composition-length bound to be proved here.

We first pass to a semisimple natural section, losing at most two trivial
factors, and then bound three sources of composition length.  A soluble
centralizer bounds the multiplicity of each absolute simple type.
Many distinct types could still occur, so we group them into orbits under
descent, the normalizer action, and duality.  The commuting-normal test
allows at most one large orbit of summands, with a bounded-dimensional
remainder.  Finally, comparing the source and ambient field shifts along
the same two-edge shortcut bounds the number of types in each orbit.
This last step makes the length bound independent of $f$.

We use Lemma~\ref{grw:coordinate-normal-tests}, with its tests
\textup{(\ref{grw:normal-soluble})} and
\textup{(\ref{grw:normal-commuting})}.
Here $P=P_{ij}$, $M=M_{ij}$, and
$F_j=\operatorname{core}_{z_j}(D)$ is the original higher core.
For an $M$-invariant $H\le T_j$, write
$\widehat H\le N_j/F_j$ for its transported product.
A killed product lies in $D/F_j$, and hence $H\le P$;
the converse is not asserted.
For one product the enlargement is
$(z_i/F_j)\widehat H$; for two it is
$\langle z_i/F_j,\widehat H_1,\widehat H_2\rangle$.
Each application checks properness of this enlargement and
normality of every tested product in that same enlargement.
The pairs used in this section commute exactly in $N_j/F_j$,
so \textup{(\ref{grw:normal-commuting})} uses $\overline K=1$.
Their full inverse images commute modulo $F_j$, which is
normal in the actual enlargement and contained in $D$.
The transported active component $\widehat E$ is not killed,
by the component-survival conclusion of
Proposition~\ref{lift:edge-lifts}.

\subsection{Semisimplicity and the nonsingular-line exception}

\begin{lemma}[Exclusion of invariant parabolics]
\label{pkt:no-parabolic}
After passing to a fixed-spacing subchain, no comparison stabilizer
$M_{ij}$ normalizes a proper parabolic subgroup of $T_j$.
\end{lemma}

\begin{proof}
This is Proposition~\ref{grw:remaining-paths}(v), with its
fixed-spacing thinning already incorporated in the retained chain.
\end{proof}

\begin{lemma}[A semisimple natural section]
\label{pkt:semisimple}
Let $V$ be the upper natural module in a retained comparison.
As a module for the lower quasisimple component, either $V$ is
semisimple, or the characteristic is two, $V$ is an even-dimensional
quadratic space, and there is an $M$-invariant nonsingular line
$R$ such that
\[
                  V^{\circ}=R^\perp/R
\]
is a semisimple symplectic module.  In the latter case the
composition length of $V$ is that of $V^{\circ}$ plus two.
The lower component and its soluble-quotient and normal-subgroup
properties pass to $V^{\circ}$.
\end{lemma}

\begin{proof}
Work first over the finite natural scalar field.  Put
\[
 S=\operatorname{soc}_E(V),\quad
 R_0=\operatorname{rad}_E(V),\quad
 U=S\cap R_0,\quad W=S+R_0.
\]
If $V$ is not semisimple, $R_0\ne0$, and a simple submodule of
$R_0$ shows that $U\ne0$.  Also $W<V$: the proper submodule $S$
is contained in a maximal submodule, which necessarily contains
$R_0$.  Socles and radicals are intrinsic, so every semilinear
normalizer preserves these constructions.

For a linear group without graph operations, $U$ gives an
$M$-normalized parabolic, contrary to
Lemma~\ref{pkt:no-parabolic}.  Suppose that $M$ also contains
duality operations.  The side-preserving subgroup has index two.
A duality acts on subspaces by a semilinear correlation, and its
transport of a module to the dual exchanges the socle with the
annihilator of the radical, and the radical with the annihilator
of the socle.  Consequently it exchanges $U$ and $W$ as members
of a flag.  Thus the flag $U\le W$ is preserved by all of $M$;
if $U=W$ it is a one-member flag.  Both members are nonzero and
proper.  This again contradicts Lemma~\ref{pkt:no-parabolic}.
This argument uses correlations on subspaces and does not treat a
graph automorphism as a linear operator on $V$.

In a bilinear or Hermitian space, identify the module with its
appropriate dual using the preserved form.  Duality of socles
and radicals gives
\[
                       R_0^\perp=S,\qquad S^\perp=R_0.
\]
Hence $U$ is totally isotropic.  This supplies a forbidden
parabolic in the symplectic, unitary, and odd-characteristic
orthogonal cases.

It remains to consider a quadratic form $Q_V$ in characteristic
two.  On a totally polar-isotropic space $U$ its restriction has
the form $\ell^2$, for a linear functional $\ell$, since the
finite field is perfect.  Its zero space is an $M$-invariant
singular subspace.  Unless $U$ is a nonsingular line, this zero
space is nonzero and contradicts
Lemma~\ref{pkt:no-parabolic}.  We are therefore left with the
nonsingular line $R=U$.

The space $R^\perp/R$ has a nondegenerate alternating form.
If its $E$-module were not semisimple, the argument with socle
and radical would give a nonzero $M$-invariant totally isotropic
subspace $A$ of that quotient.  Its inverse image $B$ in
$R^\perp$ is totally polar-isotropic and contains $R$.  On $B$
the quadratic form is the square of a linear functional nonzero
on $R$.  Its zero hyperplane is therefore a singular complement
to $R$, projecting isomorphically onto $A$.  This complement is
$M$-invariant and again gives a forbidden parabolic in the
original quadratic space.  The quotient is semisimple.

The flag $0<R<R^\perp<V$ adds precisely two one-dimensional
composition factors to that quotient.  Both are trivial for $E$
because $E$ is perfect.  By Lemma~\ref{lift:nonsingular-line},
the inner stabilizer has the canonical finite-group isomorphism
\[
 \theta:\Omega(V,Q_V)_R\xrightarrow{\sim}
                    \operatorname{Sp}(V^\circ).
\]
Every $m\in M$ induces a semilinear similarity $m^\circ$ on
$V^\circ$.  Evaluation on quotient vectors gives
\[
 \theta(mhm^{-1})=
       m^\circ\theta(h)(m^\circ)^{-1}
       \qquad(h\in\Omega(V,Q_V)_R).
\]
Consequently inverse images under $\theta$ preserve normalization
by the actual $M$, inclusions, and exact commutation of inner
subgroups.  Since $P=M\cap T_j$ preserves $R$, it lies in this
inner stabilizer.  Write $P^\circ=\theta(P)$ and
$E^\circ=\theta(E)$; then $P^\circ/E^\circ\cong P/E$.
Every candidate section subgroup is pulled back through
$\theta$ before applying a normal test.  Every lifted enlargement
lies in the proper original line stabilizer, even if its
symplectic image is the whole symplectic group.  Transport over
the original coordinates uses the same normal higher core
$F_j$ in \textup{(\ref{grw:normal-soluble})} and \textup{(\ref{grw:normal-commuting})}.  There is no additional
finite kernel in this inner-group isomorphism.

Finally, extending the finite scalar field to its algebraic
closure preserves semisimplicity here.  Indeed the semisimple
image algebra of the finite group is a product of matrix
algebras over finite fields; those fields are separable, and
scalar extension preserves semisimplicity of that algebra.
\end{proof}

\subsection{Multiplicities and packets}

In this subsection $V$ denotes the semisimple space supplied by
Lemma~\ref{pkt:semisimple}; the two removed factors will be added
back at the end.  All centralizers and groups on summands are
taken in the indicated finite classical realization and then in
the inner projective group.  Passing through the finite central
covers does not change a solubility assertion.

\begin{lemma}[Soluble centralizer]
\label{pkt:centralizer}
The centralizer of the lower quasisimple component in the inner
classical group on $V$ is soluble.
\end{lemma}

\begin{proof}
Write $C=C_{T_j}(E)$ in the original natural realization.
The group $MC$ normalizes $E$ and cannot contain $T_j$: the
proper nontrivial $E$ cannot be normal in the simple ambient
group.  Hence the full preimage of
$\overline U=\langle z_i/F_j,\widehat E,\widehat C\rangle$
is proper in $z_j$.  The two products $\widehat E,\widehat C$
are normal in $\overline U$ and commute exactly.  Their full
preimages commute modulo $F_j\trianglelefteq U$, with
$F_j\le D$.  The active product $\widehat E$ is not killed,
so \textup{(\ref{grw:normal-commuting})} kills $\widehat C$ and forces $C\le P$.
Since
\[
 C\cap E\le Z(E),\qquad C/(C\cap E)\hookrightarrow P/E,
\]
$C$ is soluble.

For the symplectic section of Lemma~\ref{pkt:semisimple}, put
$E^\circ=\theta(E)$ and take the inverse image under $\theta$
of $C_{\operatorname{Sp}(V^\circ)}(E^\circ)$.
The $M$-equivariance and injectivity of $\theta$ show that this
inner inverse image and $E$ commute exactly and are normal in
their enlargement with $M$.  Their transported products again
commute in $N_j/F_j$, and their original preimages commute
modulo the normal subgroup $F_j\le D$.  The enlargement lies
in the proper original nonsingular-line stabilizer even if
its symplectic image is the full symplectic group.
Condition \textup{(\ref{grw:normal-commuting})} and the same quotient calculation prove
that the symplectic centralizer is soluble.
\end{proof}

\begin{lemma}[An absolute multiplicity bound]
\label{pkt:multiplicity}
Every absolute irreducible constituent of $V$ has multiplicity at
most $15$.
\end{lemma}

\begin{proof}
We give the linear-algebra details needed for a bound independent
of the fields and constituent dimensions.  Let $U$ be a simple
module over the finite natural field $F$.  Its endomorphism
division ring is a finite field $F'$.  On an isotypic summand
$U^m$ the commuting endomorphism algebra is
$\operatorname{Mat}_m(F')$.  After scalar extension, $U$ splits
into distinct absolute simples, each once; thus $m$ is exactly
their common absolute multiplicity.

In the linear case the centralizer contains
$\operatorname{GL}_m(F')$.  Its subgroup
$\operatorname{SL}_m(F')$ has determinant one on the whole
$F$-space: its determinant there is a power of the norm of its
$F'$-determinant.  Thus its derived subgroup survives the
ambient special-linear condition.

For a form group, pair the simple isotypic summands by duality.
A pair of distinct dual summands again supplies
$\operatorname{GL}_m(F')$, acting on the second summand by the
inverse adjoint.  On a self-paired summand the adjoint operation
on $\operatorname{Mat}_m(F')$ is an involution.  If it acts
nontrivially on $F'$, the commuting isometries are a unitary
group on an $m$-dimensional multiplicity space.  If it fixes
$F'$, they preserve a nonsingular symmetric or alternating
bilinear form on that multiplicity space.  These descriptions
follow by writing the invariant pairing as a matrix of pairings
between the $m$ copies of $U$; Schur's Lemma makes each matrix
entry an element of $F'$, and the symmetry of the original form
is precisely the adjoint condition on that matrix.

There is one useful even-characteristic qualification.  A
nonsingular symmetric multiplicity form need not be alternating.
Over a perfect finite field its diagonal function is the square
of a linear functional.  The vector representing that functional
allows one to split off a nonsingular subspace of dimension one
or two; its perpendicular is alternating and nonsingular.
Thus a multiplicity space of dimension $m$ still contains a
nonsingular alternating subspace of dimension at least $m-2$.

If $m\ge16$, every group just described contains a subgroup
isomorphic to $\operatorname{SL}_3(F'')$ for a suitable finite
field $F''$.  In a general linear group use a three-dimensional
block.  In a symplectic, orthogonal, or unitary multiplicity
space choose paired totally isotropic three-spaces and let
$g\in\operatorname{SL}_3(F'')$ act as $g$ and its inverse
adjoint, fixing their perpendicular complement.  Finite-field
forms of dimension at least $16$ have such paired spaces;
the anisotropic part has dimension at most two.  The preceding
paragraph handles the nonalternating symmetric form in
characteristic two in the same way.

These embedded groups are perfect and have nonabelian simple
projective quotients.  Being perfect, their determinants,
similitude multipliers, spinor norms where applicable, and
Dickson invariants in even-characteristic orthogonal groups
are trivial.  Therefore
they lie in the actual inner classical group, rather than
merely its full similitude group.  A scalar quotient cannot
make their images soluble.

In an orthogonal target of characteristic two, preserving the
polar form on a nontrivial isotypic summand also preserves the
quadratic form.  Indeed two $E$-invariant quadratic refinements
of that polar form differ by $\ell^2$ for an $E$-invariant linear
functional $\ell$.  On a sum of nontrivial simple $E$-modules
there is no such nonzero functional.  A commuting polar
isometry consequently preserves the given refinement.  On the
trivial isotypic space use its actual quadratic isometry group;
the same paired-three-space construction applies.  The trivial
space is nondegenerate, since it is orthogonal to every
nontrivial isotypic summand.

We have proved that $m\ge16$ would put a nonsoluble subgroup in
the centralizer.  Lemma~\ref{pkt:centralizer} excludes this.
\end{proof}

Define a \emph{packet} to be an orbit of absolute simple types
under finite-field descent, the action induced by $M$, and the
duality pairing when a form or an ambient correlation is present,
intersected with the set of types occurring in $V$.
The sum belonging to a packet is defined over the finite natural
field.  In a form group it is nondegenerate, because all dual
partners belong to the same packet.  The packet decomposition
and every union of its parts are normalized by $M$.
For a linear graph automorphism this last statement refers to
the groups on the parts: the induced adjoint anti-automorphism
permutes the central isotypic idempotents and takes the full
linear group on a part to the full linear group on its paired
part.  It therefore normalizes the product groups used below,
although its action on subspaces is a correlation.

\begin{lemma}[At most one large packet]
\label{pkt:one-packet}
Two disjoint unions of packets on which $E$ acts nontrivially
cannot both have dimension at least $16$.  Consequently either
one packet accounts for all but at most $15$ dimensions, or
the sum of all nontrivial packets has dimension at most $45$.
\end{lemma}

\begin{proof}
Let $V_1,V_2$ be two such unions, and put all remaining parts
in $V_3$.  Let $H_1,H_2$ be the derived groups of the full
linear or isometry groups on $V_1,V_2$, acting trivially on the
other parts.  They commute, are normalized by $M$, and are
normal in the enlargement generated by them and $M$.
Its inner part preserves the proper decomposition and cannot
contain $T_j$.  Thus
$\overline U=\langle z_i/F_j,\widehat H_1,\widehat H_2\rangle$
has proper full preimage in $z_j$.  Both $H_1,H_2$ are
nonsoluble: the explicit
$\operatorname{SL}_3$ construction in the proof of
Lemma~\ref{pkt:multiplicity} lies in each derived group and in
the ambient inner group.

The two transported products are normal in this $\overline U$
and commute exactly; their original preimages commute modulo
$F_j\trianglelefteq U$, with $F_j\le D$.  Condition
\textup{(\ref{grw:normal-commuting})} kills one product and forces, say, $H_1\le P$.
The kernel
of the nontrivial action of $E$ on $V_2$ is central in $E$,
so $H_1\cap E\le Z(E)$.  Hence
\[
                 H_1/(H_1\cap E)\hookrightarrow P/E
\]
would make $H_1$ soluble, a contradiction.  On the symplectic
section, take inverse images under the finite $M$-equivariant
isomorphism $\theta$ from Lemma~\ref{pkt:semisimple}.
They commute exactly and their enlargement lies in the proper
original line stabilizer, so the same application uses the
original higher core $F_j$, as in Lemma~\ref{pkt:centralizer}.

Suppose that a nontrivial packet has dimension at least $16$.
If another nontrivial packet occurs, apply the first assertion
to the large packet and its entire complement, including the
trivial packet.  The complement has dimension at most $15$.
If no other nontrivial packet occurs, the complement is the
trivial isotypic space, whose dimension is at most $15$ by
Lemma~\ref{pkt:multiplicity}.
If no packet has dimension $16$, but their total dimension
were at least $46$, accumulate packets until their dimension
first reaches $16$.  The accumulated dimension is at most
$30$, leaving at least $16$ dimensions in the other union.
This contradicts what was just proved.  Thus in this case the
total is at most $45$.
\end{proof}

\subsection{Field shifts and the size of a packet}

The endpoint field agreement is useful only together with the
normalizer action on both edges of a shortcut.  The quantity
to control is the difference between the source field action
and the scalar action on the upper natural space.

\begin{lemma}[Field-shift synchronization]
\label{pkt:field-shifts}
Let the endpoint exponent be $f$, and let the intermediate
exponent on a remaining shortcut be
\[
                f'=\frac uv f,\qquad (u,v)=1.
\]
Put $h=\gcd(f,f')=f/v$.  If an element of $M$ induces source
field exponent $a$ and upper scalar exponent $s$, then
\begin{equation}
                   s-a\equiv0\pmod h,\qquad
                   v(s-a)\equiv0\pmod f.
\label{pkt:shift-congruence}
\end{equation}
Every packet of absolute
constituent types has size at most $2v$.
\end{lemma}

\begin{proof}
We work on simply connected covers and use Frobenius positions
modulo the group parameter, rather than the degree of its
natural scalar field.  Write a classical finite group with
parameter $p^r$ as the fixed points of $\tau F_p^r$, where
$\tau$ is a diagram operation.  Steinberg's Restriction and
Tensor Product Theorems give the unique parametrization
\[
 \lambda=\lambda_0+p\lambda_1+\cdots+p^{r-1}\lambda_{r-1},
 \qquad\lambda_i\text{ $p$-restricted},
\]
of its absolute simple modules; see
\cite[Theorems~41 and~43]{SteinbergNotes} and
\cite{Steinberg1963}.  For a nontrivial simple module define
\[
 \operatorname{supp}_r(\lambda)
       =\{i\in\mathbb Z/r\mathbb Z:\lambda_i\ne0\}.
\]
For an arbitrary module take the union of these supports over
its nontrivial composition factors.  A field twist by $F_p^b$
translates support by $b$.  Diagram operations and
contragredience preserve positions.  When a position wraps
modulo $r$, its restricted weight is transformed by
$\tau^{-1}$, which does not change the position.

Consider one edge, with source and target exponents $r,t$, and
let $d=\gcd(r,t)$.  Suppose that the nonempty union $\mathcal E$
of source digit supports in the target natural module lies in
one residue class modulo $d$.  If a normalizer has target
scalar exponent $b$ and source field exponent $a_0$, its
transport of a composition factor has type
\[
                  (L(\lambda)\circ\gamma^{-1})^{[b-a_0]},
\]
where $\gamma$ records the source diagram operation.  This
formula follows directly from
$x(mw)=m(\alpha_m^{-1}(x)w)$ and the $p^b$-semilinearity of
$m$.  It applies to composition factors by transporting a
composition series; semisimplicity on an individual edge is
not required.  An ambient correlation adds contragredience.
Inner-diagonal automorphisms do not affect the absolute module
type, because they are algebraic inner conjugations on the
cover.  Normalization therefore implies
\[
                  \mathcal E+(b-a_0)=\mathcal E.
\]
Since $\mathcal E$ is nonempty and contained in a single
residue class modulo $d$, we obtain
\begin{equation}
                          b-a_0\equiv0\pmod d.
\label{pkt:one-edge-shift}
\end{equation}

Assume for the moment that this support condition holds for both
edges of the shortcut.  Their source and target exponents are $f,f'$
in opposite orders, so their common divisor $d$ is the same integer $h$.
For an element of $M$, denote its induced intermediate field
exponent by $a_1$.  Its existence on the same intermediate
component is the normalization conclusion of
Proposition~\ref{lift:edge-lifts}.  The natural realization of that
automorphism has scalar exponent $a_1$, with any graph part
recorded separately.  Applying
\eqref{pkt:one-edge-shift} to the two edges gives
\[
                 a_1-a\equiv0\pmod h,\qquad
                 s-a_1\equiv0\pmod h.
\]
Their sum proves \eqref{pkt:shift-congruence}.

To bound a packet, let $\Phi$ denote the operation of field twisting on
absolute source types, and let $\Gamma$ be the diagram group.
The rank assumption gives $|\Gamma|\le2$.  These operations
commute, and $\Phi^f=\tau^{-1}$.  Every stabilizer action lies
in $\langle\Phi^h,\Gamma\rangle$.  If $\tau=1$, that group has
order at most $2f/h$.  If $\tau\ne1$, the diagram involution
already belongs to $\langle\Phi^h\rangle$, whose order is at
most $2f/h$.  Thus in either case the order is at most $2v$.
Descent over the upper natural field has exponent $f$ or $2f$
and belongs to the same group.  Duality is the opposition
diagram operation, also already included.  The packet bound
follows.

It remains to verify the support hypothesis for every remaining edge of
Proposition~\ref{grw:remaining-paths} and
Proposition~\ref{lift:edge-lifts}.
\begin{enumerate}[label=\textup{(\alph*)}]
\item In a nonlinear edge the natural module is a Frobenius
      twist of a nonzero restricted module.  Its support is
      a singleton.  The starting Frobenius position may be
      arbitrary; it cancels in
      \eqref{pkt:one-edge-shift}.
\item A block, form, or subfield edge has natural or dual-natural
      nontrivial composition factors in one Frobenius position,
      allowing repetitions and trivial factors.  A tensor edge
      with one active component has the same property: its
      other tensor factors are trivial for that perfect
      component, so they supply repetitions.  The
      characteristic-two nonsingular-line representation has
      these composition factors as well, with the additional
      trivial factors already described.
\item For restriction of scalars, the source constituents are
      a direct sum of the Galois conjugates of the natural
      module.  To distinguish group parameters from natural
      scalar fields, put $\epsilon_s,\epsilon_t=2$ for a
      unitary source or target, respectively, and $1$
      otherwise.  A scalar restriction of degree $b_{\rm sc}$
      satisfies
      \[
          \epsilon_s r=b_{\rm sc}\epsilon_t t.
      \]
      Its constituent positions are
      $t_0+j\epsilon_t t$, so they all lie in one residue
      class modulo $d=\gcd(r,t)$.
      For example, restricting the natural matrix field of a
      unitary component with parameter $p^r$ to a symplectic
      target with the same parameter has scalar degree two.
      Over the algebraic closure its module is
      \[
                         N\oplus N^{[r]}
                         \cong N\oplus N^*.
      \]
      Both positions are zero modulo $r$; this repetition
      causes no tensor carry because the operation is a
      direct sum.
\item For a retained natural twisted-tensor edge, let
      $b_{\rm cyc}$ be the prime tensor-cycle length.
      Lemma~\ref{cls:twisted-tensor} gives the exact usual-field
      relation
      \[
                           r=b_{\rm cyc}t.
      \]
      In the finite realization supplied by
      Proposition~\ref{lift:edge-lifts}, the cyclic fixed-point
      construction gives the following actual finite-source
      modules, up to a common initial Frobenius twist
      \cite[Proposition~2.8(ii) and Section~4(5), p.~449]
      {LiebeckSeitzClassical1998}.
      Here $W$ is the source natural module in each case.

      \emph{Linear, symplectic, and orthogonal targets.}
      The module is
      \[
                   \bigotimes_{j=0}^{b_{\rm cyc}-1}W_j^{[jt]},
      \]
      where each $W_j$ is $W$ or $W^*$ as prescribed by the
      diagram operation.  Its positions are $jt$ for
      $0\le j<b_{\rm cyc}$; diagram changes affect the
      weights rather than these positions.

      \emph{Unitary target, odd cycle.}
      The source is $\operatorname{SU}_m(p^r)$ and the module is
      \[
                   \bigotimes_{j=0}^{b_{\rm cyc}-1}W^{[2jt]}.
      \]
      Indeed, the cycle initially gives $W^{[jt]}$ for even
      $j$ and $(W^*)^{[jt]}$ for odd $j$.  On this finite source,
      $W^*\cong W^{[r]}$.  The integers $j$ for even $j$ and
      $j+b_{\rm cyc}$ for odd $j$ are precisely
      $0,2,\ldots,2b_{\rm cyc}-2$, giving the displayed module.
      Its positions $2jt$ are distinct modulo $r$ because
      $b_{\rm cyc}$ is odd.

      \emph{Unitary target, two-cycle.}
      Since the cycle length is prime, this is the only even
      case.  The source is $\operatorname{SL}_m(p^{2t})$ and
      its module is
      \[
                             W\otimes(W^*)^{[t]},
      \]
      with distinct positions $0,t$ modulo $r=2t$.

      Thus in every case the positions are distinct modulo
      $r=b_{\rm cyc}t$ and lie in one residue class modulo
      $d=\gcd(r,t)=t$.  A common initial twist adds the same
      $t_0$ to every position and preserves both assertions.
      Steinberg's Tensor Product Theorem therefore gives the
      required support condition with no coincident-digit
      carries.  This distinctness comes from the exact finite
      cyclic construction.

\end{enumerate}
Changes between the type $B$ and type $C$ conventions in
characteristic two do not alter the numerical field exponent.
The special isogenies are defined over $\mathbb F_2$, commute
with Frobenius, and have Frobenius as their composites
\cite[Section~I.3.1]{DowdSin1996}.  Thus each edge can use the
convention in which its displayed module has the indicated
support, while the intermediate field action has the same
exponent in either convention.
\end{proof}

\begin{theorem}[Uniform absolute composition length]
\label{pkt:bounded-length}
In the setting of Proposition~\ref{grw:remaining-paths}, let $B$ be
an absolute bound for the denominators of the intermediate
field ratios $f'/f$ on the remaining shortcuts.  After a
further fixed-spacing thinning, every upper natural module,
restricted to a lower active quasisimple component, has at most
\[
                              L=30B+62
\]
absolute composition factors, counted with multiplicity.
The bound uses neither natural-module irreducibility between
endpoints nor an irreducible-triples theorem.
\end{theorem}

\begin{proof}
Apply Lemma~\ref{pkt:semisimple}.  On its semisimple natural
section, Lemma~\ref{pkt:multiplicity} bounds each absolute
multiplicity by $15$, and Lemma~\ref{pkt:field-shifts} bounds
each packet by $2B$ absolute types.

If there is a large nontrivial packet, it contributes at most
$30B$ factors.  All other nontrivial packets have total
dimension at most $15$ by Lemma~\ref{pkt:one-packet}.
If no packet is large, their total dimension is at most $45$.
In either case the number of nontrivial composition factors
is at most $30B+45$.  The trivial isotypic space has dimension
at most $15$, again by Lemma~\ref{pkt:multiplicity}.  Thus the
semisimple section has length at most $30B+60$.  Adding the
possible two factors from the nonsingular-line case gives the
displayed bound.

All constants came from the remaining-edge bounds, explicit
finite-dimensional linear algebra, and the common cyclic
congruence \eqref{pkt:shift-congruence}.  They are independent
of field exponent, rank, and chain length.
\end{proof}
